1. Relationship Between the Two Concepts

True repayment fraction** \(\rho_t\):  
  The share of newly issued claims that are actually extinguished against real output in period \(t\) (primarily via primary surplus).

Fraction Never Repaid in Real Purchasing Power**:  
  The cumulative complement of true repayment over a multi-period horizon:

\[
1 - \frac{\sum pb_t}{\lvert d_T - d_0 \rvert}
\]

This cumulative fraction is the historical aggregate expression of low average \(\rho_t\).

In short:
\[
\text{Average }\rho \approx 1 - \text{(Fraction Never Repaid)}
\]

2. Empirical Content from 1900s–2000s

Historical decompositions across the 20th and early 21st centuries show:

True repayment via sustained primary surpluses was the exception rather than the rule.
In most periods, especially after the 1970s and particularly post-1980 / post-2008, \(\rho_t\) remained low.
A large share of claims was handled through rollover, growth, inflation, financial repression, or outright restructuring/default-like outcomes.
This pattern effectively normalized non-repayment (or partial, delayed, and incomplete repayment) as standard practice.

The data therefore support treating the high “never repaid in real purchasing power” fraction as the dominant empirical reality, which in turn implies a structurally low average \(\rho_t\).

3. Model Implication

Under this reading:

\[
\rho_t \text{ is low on average (typically 0.10–0.25, often lower)}
\]

and the complementary high never-repaid fraction becomes the primary driver of residual accumulation and the associated growth of the unextinguished claim stock.

This is now the working empirical stance: low true repayment rates, validated by 20th- and early 21st-century experience, form the analytical basis for residual growth.
